\documentclass[sigconf,nonacm]{acmart}
\settopmatter{printacmref=false}
\setcopyright{none}
\renewcommand\footnotetextcopyrightpermission[1]{}
\pagestyle{plain}
\raggedbottom

\acmConference[Manuscript Draft]{Unicode Readiness Study}{2026}{}
\acmYear{2026}
\acmISBN{}
\acmDOI{}

\title{Unicode Readiness for Nwagu Aneke: A Gap Matrix for Character Evidence and Community Review}
\author{Anonymous Author(s)}
\affiliation{\institution{Independent research manuscript}\country{Nigeria}}
\email{anonymous@example.invalid}

\begin{abstract}
Underdocumented scripts often enter digital-humanities and standards discussion
before their evidence is strong enough for character encoding. This paper
examines that intermediate condition through Nwagu Aneke, an Igbo writing
system described in the scholarly and public script-status literature. The
source base used here preserves a 26 by 8 table of vowel/modifier combinations,
yielding 208 source-observed records, while treating a 27/216 f/v expansion as
derived. A twelve-item Unicode-readiness matrix evaluates script identity,
repertoire, character-glyph distinction, representative glyphs, character
names, encoding model, directionality, punctuation and numbers, usage examples,
community authority, implementation evidence, and proposal dossier status. The
matrix finds one present requirement, four partial requirements, three blocked
requirements, and four missing requirements. The result supports an evidence
audit of readiness conditions while withholding encoding, repertoire, glyph,
and public data-release conclusions until source transcription review, rights
review, and community authority review are completed.
\end{abstract}

\keywords{Nwagu Aneke, Igbo writing systems, Unicode readiness, digital humanities, script documentation}

\begin{document}
\maketitle

\section{Introduction}
Digital work on underdocumented scripts faces a recurring timing problem. A
script may be documented well enough to be discussed, compared, and modeled, but
too incompletely documented to support encoding, public glyph charts, or
software implementation. If that intermediate state is handled poorly, the
researcher has only two unsatisfactory choices: either avoid the script until a
complete edition exists, or write as though partial evidence already supports a
standards proposal. Neither choice is adequate for cultural heritage, digital
humanities, or language technology.

Nwagu Aneke makes the timing problem concrete. Azuonye's account describes the
script as an Igbo writing system with potential significance for alternative
literacy \cite{azuonye1992}. Public script-status resources identify the script
and connect it to Unicode-status discussion \cite{omniglotNwagu,scriptSourceNwagu}.
At the same time, the evidence currently available for computational and
standards-facing work remains incomplete. The working source base used in this
paper is a 26 by 8 table of vowel/modifier combinations, for 208
source-observed records. A 27/216 representation created by splitting f/v is
treated as a derived analytical layer. That distinction matters because a useful
derived representation can become misleading if it is later described as
source-observed.

This paper asks whether the intermediate state can be made analytically useful
without making it look more mature than it is. The research question is: what
Unicode-readiness requirements are present, partial, blocked, or missing for
Nwagu Aneke under the present evidence base? The hypothesis is that a
requirement-by-requirement matrix can produce a useful negative result: it can
show why the script is suitable for readiness analysis while still withholding
encoding, repertoire, glyph, and authority conclusions.

The contribution is a matrix and its interpretation. Twelve requirements are
evaluated: script identity, repertoire inventory, character-glyph distinction,
representative glyphs, names list, encoding model, directionality, punctuation
and numbers, usage examples, community authority, implementation evidence, and
proposal dossier. The matrix finds one present requirement, four partial
requirements, three blocked requirements, and four missing requirements. That
distribution is the paper's central result. It is not a failure to produce an
encoding proposal. It is evidence that the responsible next work is source
review, rights review, authority review, and character-evidence collection.

\section{Related Work}
The source anchor is Azuonye's study of the Nwagu Aneke Igbo script
\cite{azuonye1992}. Public descriptions and script-status records provide
secondary context for the script's visibility and current encoding status
\cite{omniglotNwagu,scriptSourceNwagu,sirisNwaguProposal}. The Unicode
Consortium's African-script update and current Unicode documentation provide
the standards context in which unencoded and proposal-stage scripts are
discussed \cite{unicodeAfricanScripts2023,unicode17CoreSpec}. Unicode proposal
guidance and roadmaps clarify that a script's cultural or historical interest
does not by itself settle the evidence required for encoding
\cite{unicodeProposals,unicodeRoadmaps,unicodePipeline2026,seiProposalTips}.

The paper also draws on digital-humanities and cultural-heritage infrastructure.
TEI offers mechanisms for describing characters, glyphs, and writing modes
\cite{teiP5Glyphs}. IIIF and the Web Annotation Data Model support source
presentation and annotation of digitized cultural materials
\cite{iiifPresentation3,w3cAnnotationModel,w3cAnnotationVocab}. PROV-O,
RO-Crate, DataCite, and CIDOC CRM provide vocabularies for provenance,
research-object packaging, citation, and cultural-heritage modeling
\cite{w3cProvO,rocrate12,datacite45,cidocCrm}. FAIR and CARE principles frame
the difference between technical reusability and responsible stewardship,
especially where community authority and cultural data are involved
\cite{wilkinson2016fair,carroll2020care}.

Unicode readiness is best understood as a layered evidence problem rather than
as a simple question of whether a script is interesting enough to encode. The
Script Encoding Working Group guidelines foreground the distinction between
characters and glyphs, the need to show usage, stability, and interchange need,
and the practical requirements of proposal documents, properties, summary
forms, fonts, and licensing \cite{unicodeProposals}. Roadmap pages make a
related distinction visible: some scripts are already published, some have
formal or mature proposals, and some remain exploratory or provisional
\cite{unicodeRoadmaps}. For a script such as Nwagu Aneke, these distinctions
matter because an exploratory source table, a digital-humanities model, and a
standards submission are different scholarly objects. Treating them as one
object would make a partial evidence base appear more settled than it is.

The character-glyph distinction is central to the present study. Unicode
encoding is concerned with abstract characters and their use in plain-text
interchange, while glyphs are visual realizations that may vary across hands,
fonts, manuscripts, or editorial displays \cite{unicode17CoreSpec}. This
distinction becomes difficult in historical and underdocumented scripts because
the surviving evidence may consist of charts, pedagogical displays, isolated
examples, or manuscript images rather than a settled typographic tradition. A
row in a chart may be a character, a class of glyphs, an editorial label, or a
convenient representation of a larger orthographic pattern. The present matrix
therefore does not move directly from a visible table to a repertoire claim. It
asks which intermediate conditions have been satisfied before such a move would
be responsible.

Those infrastructures are relevant because Unicode readiness is not only a
standards problem. Before a proposal exists, a research community may still need
to describe source witnesses, record uncertainty, compare transcriptions, cite
metadata, and share derived representations. TEI can represent a distinction
between a glyph occurrence and a normalized character-like unit. IIIF and Web
Annotation can identify regions of source images without forcing those regions
to become public glyphs. PROV-O and RO-Crate can describe how a matrix was
derived from sources. CIDOC CRM can place signs, documents, and events in a
cultural-heritage model. These systems do not answer the encoding question by
themselves. They provide a way to preserve evidence discipline before encoding
language becomes appropriate.

Finally, the problem is related to evidence preservation in computational
systems. Scientific claim verification and citation-grounded generation
research show how fluent text can misstate what cited evidence supports
\cite{wadden2020scifact,gao2023alce}. Tokenization research shows that
computational representations can be useful while also encoding design choices
that affect later systems \cite{sennrich2016bpe,kudo2018sentencepiece,bostrom2020bytepair,rust2021tokenization}.
For Nwagu Aneke, the analogous risk is that a source table, a derived f/v
expansion, and a future character model may be collapsed into one apparent
inventory. Unicode readiness therefore requires evidence typing as well as
evidence collection.

The analogy to tokenization is methodological rather than topical. A tokenizer
does not merely count units; it defines a representational layer whose choices
affect downstream models. A script-readiness matrix likewise defines what kind
of unit is being counted. In the Nwagu Aneke case, 208 source-observed records
and 216 derived analytical records are both useful numbers, but they answer
different questions. The former describes the current source table used here;
the latter describes an expansion produced by splitting f/v for analysis. A
standards-facing character inventory would require a third object: a reviewed
set of proposed abstract characters with names, properties, representative
glyphs, usage evidence, and authority review. Much of the present paper is
therefore concerned with preventing movement between these three objects before
the evidence supports it.

\section{Materials and Evidence}
The material examined here is a text-only evidence base for Nwagu Aneke. The
accepted source layer is a 26 by 8 arrangement of consonant or row labels against
vowel/modifier columns, producing 208 source-observed records. The derived f/v
layer expands that structure to 27/216 for analytical purposes, but the
expansion is not treated as primary evidence. This separation is retained
throughout the matrix.

The matrix uses three evidence categories. The first category is source
evidence: the specific records that support the 26 by 8 table and the textual
claims that identify Nwagu Aneke as an Igbo writing system. The second category
is standards evidence: public Unicode, Script Encoding Working Group, and
Script Encoding Initiative materials that describe what a proposal would need to
contain. The third category is stewardship evidence: source-review, rights, and
authority conditions that determine whether examples, glyphs, and manuscript
data can responsibly be made public. A row is treated as strong only when the
relevant category is present. For example, script identity can be present with
source and public metadata evidence; representative glyphs require source
images and rights review; community authority requires a different kind of
evidence altogether.

The text-only scope is important. This paper does not reproduce source images,
glyph crops, manuscript examples, or a font. It also does not introduce a
private-use encoding, a keyboard, or a proposed block allocation. The absence of
these artifacts is not hidden. It is recorded as part of the result. In a
readiness audit, evidence that is absent or unavailable is just as important as
evidence that is present, because absent evidence determines the ceiling of the
public claim.

The readiness matrix is summarized in Table~\ref{tab:readiness}. A status of
``present'' means the evidence is adequate for the matrix row at the audit
stage. A status of ``partial'' means relevant information exists, but further
review or formalization is needed. A status of ``blocked'' means the row depends
on source access, rights clearance, or authority decisions that cannot be
substituted with formatting or inference. A status of ``missing'' means the
current evidence base lacks the relevant material.

\begin{table*}[t]
\caption{Unicode-readiness matrix for Nwagu Aneke under the current evidence base}
\label{tab:readiness}
\begin{tabular}{p{0.18\textwidth}p{0.11\textwidth}p{0.62\textwidth}}
\toprule
Requirement & Status & Evidence or gap \\
\midrule
Script identity & Present & Nwagu Aneke is identified in Azuonye's account and public script metadata. \\
Repertoire inventory & Partial & The 26 by 8 source table exists; glyph-level repertoire review remains incomplete. \\
Character-glyph distinction & Partial & A working distinction can be stated, but expert encoding review is still needed. \\
Representative glyphs & Blocked & Public image rights and selector-level review are unresolved. \\
Names list & Partial & Row and vowel readings exist; Unicode character names have not been drafted. \\
Encoding model & Missing & No formal block, ordering, property, or character model is available. \\
Directionality & Partial & Public descriptions suggest left-to-right writing; source confirmation remains necessary. \\
Punctuation and numbers & Missing & No reviewed evidence is available in the current package. \\
Usage examples & Blocked & Manuscript corpus examples are not rights-cleared for public use. \\
Community authority & Blocked & Authority and consultation records remain incomplete. \\
Implementation evidence & Missing & No font, keyboard, or interoperable encoding implementation is available. \\
Proposal dossier & Missing & No proposal dossier exists under the standards requirements considered here. \\
\bottomrule
\end{tabular}
\end{table*}

\section{Method}
The method has three steps. First, the current source and derived layers are
separated. The source-observed layer is the 26 by 8 arrangement; the f/v
expansion is recorded as derived. Second, a readiness rubric is constructed from
the kinds of evidence normally needed before unencoded scripts can be discussed
as candidates for character encoding: identity, repertoire, character-glyph
distinction, representative glyph evidence, names, model, directionality,
punctuation and numbers, usage examples, authority, implementation, and proposal
documentation. Third, each row is assigned a status and interpreted in relation
to the strongest public claim it can support.

The research question is operationalized as follows: for each readiness
requirement, what is the strongest public statement supported by the evidence
now available? The hypothesis is that the answer will be uneven across
requirements: some rows will support limited positive claims, while others will
show that the current materials are insufficient for proposal framing. The
method therefore evaluates each row twice. The first pass asks whether a
specific evidence item exists. The second pass asks whether that evidence item
is the right kind of evidence for the row. A public description of direction is
useful, for instance, but it is weaker than a reviewed set of source examples
showing writing direction across documents. A chart is useful for repertoire
discussion, but it is weaker than a reviewed character inventory.

The method deliberately avoids a binary ready/unready label. Binary labels are
too coarse for underdocumented scripts. They hide whether the problem is
technical, evidential, legal, or community-governance related. The present
matrix instead distinguishes four conditions. Present rows can be used as
starting points. Partial rows invite further documentation. Blocked rows require
human source transcription review, rights review, or authority review. Missing
rows identify evidence that has not yet been collected.

The unit of analysis is the requirement row, not the script as a whole. This is
why the matrix can say that script identity is present while implementation
evidence is missing. It can also say that repertoire inventory is partial while
representative glyphs are blocked. A whole-script label would flatten those
differences. Row-level analysis preserves them and makes the next task more
precise. If the blocked rows remain blocked, no amount of prose polish can
repair the readiness gap. If the missing rows are later supplied, the matrix can
be updated without rewriting the source-observed and derived-layer distinction.

The row statuses were assigned by four coding rules. The first rule is
specificity: each row must point to a named evidentiary object or to a named
absence. For example, the repertoire row points to the 26 by 8 source layer, and
the implementation row points to the absence of a font, keyboard, or
interoperable encoding test. The second rule is type fit: the evidence must
match the decision being made. A public description can support script identity,
but it cannot replace representative glyph evidence. A source table can support
inventory discussion, but it cannot replace a character-property model. The
third rule is reversibility: if a future reviewer obtains stronger evidence,
the row should be updatable without disturbing unrelated rows. A rights decision
about glyph crops, for instance, would affect representative glyphs and perhaps
usage examples, but it would not automatically settle punctuation or ordering.
The fourth rule is claim restraint: a row may support only the strongest claim
licensed by its weakest necessary evidence category.

Those coding rules make the matrix inspectable. A reviewer can disagree with a
row by naming which rule was applied incorrectly. If the reviewer believes that
directionality should be present rather than partial, the reviewer must show
that the existing evidence is both specific and type-fitting for directionality.
If the reviewer believes that representative glyphs should be partial rather
than blocked, the reviewer must show that source access and public-use rights
are sufficient for the glyph examples being discussed. If the reviewer believes
that implementation evidence should be partial rather than missing, the
reviewer must identify an implementation artifact that can be inspected outside
the prose. This is more useful than a general readiness score because it turns
disagreement into a row-level evidence question.

The method also distinguishes three forms of uncertainty. Transcription
uncertainty concerns whether the source table has been read correctly. Modeling
uncertainty concerns whether a source record should become a character, a glyph
variant, a sequence, or a derived analytical unit. Stewardship uncertainty
concerns whether public release of an example is authorized and appropriate.
These uncertainties interact but should not be collapsed. A perfectly
transcribed table could still lack a defensible encoding model. A technically
coherent model could still lack rights clearance for examples. A rights-cleared
example could still require community review before being used as a
representative public form. The matrix keeps those uncertainties visible by
assigning blocked and missing rows separately from partial rows.

The source/derived distinction is the main control condition. Every matrix row
is read against the rule that 26 by 8 is the current source-observed structure
and 27/216 is a derived f/v expansion. The distinction is not merely arithmetic.
It determines whether a later reader is looking at source evidence, editorial
analysis, or a hypothetical standards object. In an underdocumented script,
numbers can become persuasive even when their evidence layer is weak. The
matrix therefore uses the source/derived distinction as a guardrail for all
readiness claims.

Finally, the method treats blocked rows as first-class findings. In many
technical reports, rights, authority, and source-access questions appear as
administrative caveats after the main result. That ordering is unsuitable here.
For cultural-heritage scripts, rights and authority affect which evidence can
be inspected, which examples can be reproduced, and whose knowledge is being
represented. A blocked row is therefore part of the result, not an external
obstacle to the result. This is why the matrix can produce a useful negative
finding: it identifies exactly which prerequisites must change before the
stronger standards question can be reopened.

The method also treats negative results as results. A missing encoding model is
not simply an absence; it is a constraint on what scholars, technologists, and
standards participants should write next. A blocked representative-glyph row is
not a typographic inconvenience; it marks the point where rights and source
access matter. A missing proposal dossier prevents proposal framing even if the
script is culturally important. The matrix is useful only if these limits are
visible.

\section{Results}
The matrix assigns one requirement to the present category: script identity.
Nwagu Aneke can be identified and discussed as a historically meaningful Igbo
writing system on the basis of existing scholarship and public metadata
\cite{azuonye1992,omniglotNwagu,scriptSourceNwagu}. This supports scholarly
discussion of the script and justifies a readiness audit.

The single present row should not be underestimated. Script identity is the row
that makes the rest of the audit meaningful. Without an identifiable script,
the matrix would only be an abstract exercise. Here, the script can be named,
cited, and placed in a public script-status conversation. That is enough to
justify an evidence audit. It is not enough to settle the internal structure of
the script, but it does prevent the case from being treated as merely
hypothetical.

Four requirements are partial. Repertoire inventory is partial because the 26 by
8 table gives a structured basis for discussion, while glyph-level character
review remains incomplete. Character-glyph distinction is partial because the
paper can distinguish source records, glyph appearances, and derived analytical
forms, but an expert encoding analysis has not yet reviewed which distinctions
should become characters. Names list is partial because row and vowel readings
exist, but formal character names have not been drafted. Directionality is
partial because public descriptions indicate left-to-right writing, while
source confirmation is still needed.

The partial rows identify where the current evidence is useful but unstable.
The 26 by 8 table supports a controlled discussion of source-observed records,
yet it does not decide whether each record corresponds to a Unicode character.
The character-glyph distinction can be stated in principle, yet the distinction
has not been applied across a reviewed corpus of examples. Row and vowel
readings provide material for future names, yet Unicode names must follow
standards conventions and require review. Directionality can be described from
public accounts, yet a proposal would need stronger source examples and text
behavior analysis. These are not empty gaps; they are the places where the next
evidence collection would most efficiently increase readiness.

Three requirements are blocked. Representative glyphs are blocked because public
image rights and selector-level review are unresolved. Usage examples are
blocked because manuscript corpus examples cannot be assumed available for
public distribution. Community authority is blocked because consultation and
authority records are incomplete. These rows cannot be solved by inferential
confidence. They require decisions by people with appropriate source, rights,
and community knowledge.

The blocked rows are qualitatively different from the partial rows. They cannot
be advanced by more careful wording alone. Representative glyphs require a
decision about which visual forms can be shown and under what license. Usage
examples require source access and permissions that allow readers to inspect
the examples. Community authority requires accountable consultation, especially
because Nwagu Aneke is not only an object of technical interest but part of an
Igbo cultural and historical setting. Treating these rows as merely technical
would be a category error.

Four requirements are missing. No formal encoding model is available. No
reviewed evidence for punctuation and numbers is present. No implementation
evidence such as a font, keyboard, or interoperable encoding experiment is
available. No proposal dossier exists. These missing rows are decisive for
standards framing. A readiness discussion can proceed, but proposal language
would outrun the evidence.

The missing rows also show why the present work should remain a readiness
study. An encoding model would need decisions about repertoire structure,
ordering, properties, combining behavior if any, normalization implications, and
plain-text interchange. Punctuation and numbers would need either script-specific
evidence or an explicit argument for using existing characters. Implementation
evidence would need an appropriate font or other interoperable demonstration.
A proposal dossier would need a coherent document, forms, property data,
examples, and licensing. None of these items follows automatically from the 26
by 8 table.

The resulting distribution is therefore 1 present, 4 partial, 3 blocked, and 4
missing. The distribution shows why the strongest current contribution is an
evidence-gap analysis. The source base is strong enough to organize the problem,
but not strong enough to settle repertoire, glyph, usage, implementation, or
authority questions.

The distribution has a further implication: the order of future work matters.
The first priority is not a font, because representative glyphs and usage
examples remain blocked. The first priority is also not a proposal document,
because the encoding model and dossier are missing. The most efficient next
work is to stabilize the evidentiary base: review the source transcription,
resolve public rights questions for examples, document authority relationships,
and then decide whether the partial rows can be promoted to present. Only after
that sequence would implementation or proposal work become analytically
well-founded.

\section{Discussion}
The readiness matrix changes the form of the scholarly question. Instead of
asking whether Nwagu Aneke should be encoded, it asks what evidence would have
to exist before encoding discussion could responsibly proceed. That shift is
important. Encoding is not only a technical desire to preserve a script in a
universal character set. It is also a claim about character identity,
representative forms, usage, stability, and community interest.

This reframing is especially important for scripts whose public evidence is
fragmentary. If the first public computational representation is too confident,
later work may inherit its categories as facts. A derived f/v expansion may
become a supposed source inventory; a row label may become a character name; a
chart cell may become a glyph standard. Once those moves are embedded in
datasets, software, or teaching materials, they become harder to correct. The
matrix is therefore a preventive instrument. It records where confidence is
earned and where it is still pending.

The matrix also separates three activities that are often conflated. The first
activity is source-critical documentation: preserving what can be said from
available records. The second is digital-humanities modeling: representing
source and derived layers in TEI, annotation, provenance, or research-object
form. The third is standards proposal work: preparing evidence for character
encoding. Nwagu Aneke can support the first two activities in limited ways
before it supports the third.

This separation has practical value. A digital edition can represent the 26 by
8 source layer and label the f/v expansion as derived. A teaching or research
interface can show the difference between observed and analytical layers. A
metadata record can identify the script and cite the scholarly source. None of
those activities requires a public code chart or a completed Unicode proposal.
Conversely, a proposal would require a much stronger evidentiary base than the
present matrix supplies.

The separation also helps explain what kind of public artifact would be
appropriate at each stage. A source-critical note can report the 26 by 8
structure and the derived status of 27/216. A technical report can describe how
to store uncertainty and provenance. A future digital edition could expose
cleared examples with annotations. A future proposal could draw on those
materials only after they have been reviewed. The same evidence may travel
across these artifacts, but its claim status changes as the review context
changes. This is why the present paper treats readiness as a sequence of
evidence states rather than a single manuscript milestone.

The blocked rows are especially important for cultural-heritage practice. A
script can be historically significant and still require careful handling of
source images, manuscript examples, and authority claims. FAIR principles make
data findable and reusable; CARE principles remind researchers that reusability
alone is not enough where people, communities, and authority relationships are
involved \cite{wilkinson2016fair,carroll2020care}. The matrix therefore treats
rights and authority as part of the evidence condition rather than as an
afterthought.

For standards work, the same point appears as a licensing and process issue.
Unicode proposal guidance requires attention to contributors, potential IP
holders, fonts, proposal documents, and supporting evidence
\cite{unicodeProposals}. Those requirements are not bureaucratic decoration.
They are the mechanism by which a proposed character set becomes usable in a
global standard without creating avoidable ambiguity or rights problems. Nwagu
Aneke therefore needs both philological care and standards care. The former
asks whether the source reading is correct; the latter asks whether the
proposed character model would support public interchange.

The analysis also suggests a modest role for computational tools. Tools can
help maintain separate layers, check whether derived counts are being promoted
into source counts, and preserve provenance from table rows to later models.
They can also make absence visible by keeping blocked and missing rows in the
same table as present rows. But tools cannot create authority, licensing, or
source access. A well-designed computational workflow should therefore slow
down at exactly the points where this matrix marks a row as blocked.

The matrix can also serve as a comparison template for other underdocumented
scripts, provided that its categories are not treated as universal scores. A
different script may have strong manuscript evidence but weak implementation
evidence, or strong community usage but incomplete property analysis. The value
of the template is the discipline of separating requirement rows, evidence
types, and claim strength. For Nwagu Aneke, that discipline prevents the 208
source-observed records from being inflated into a settled repertoire. For
another script, the same discipline might prevent a font demonstration from
being mistaken for usage evidence, or a community teaching chart from being
mistaken for a stable encoding model. The comparative lesson is that readiness
is not a single ladder. It is a set of evidence relations that mature at
different rates.

For that reason, a future revision should report row movement explicitly. A row
should move from missing to partial only when a concrete artifact appears, from
partial to present only when the artifact fits the decision, and from blocked
to partial only when the blocking review has actually occurred.
This preserves reviewable evidence history.

This point matters for digital preservation. Preservation work often begins
before all standards questions are resolved. A repository can preserve a source
description, a transcription, a provenance trail, or a non-normative teaching
representation while still making clear that the material is not a character
encoding. Such early preservation is valuable when it keeps uncertainty
visible. It becomes risky only when provisional structures are exported as
settled public standards. The matrix therefore supports preservation and
restraint at the same time: it encourages documentation of what is known, while
making the missing and blocked requirements hard to ignore.

\section{Limitations}
The matrix is limited by the evidence base available for this paper. It does not
provide a public glyph corpus, a source-image edition, a manuscript corpus, or a
character-name proposal. The source-observed 26 by 8 structure may need revision
if future source transcription review identifies errors or omissions. The
derived f/v expansion may remain useful for analysis, but it should not be
treated as source-observed without additional evidence.

The matrix is also limited by its granularity. Twelve rows are enough to expose
the major readiness gaps, but they do not capture every technical detail that a
full script proposal would require. Character properties, ordering behavior,
line breaking, shaping behavior, normalization, representative forms, and font
engineering would each require more detailed treatment. The row labeled
``encoding model'' therefore compresses several future technical questions. It
is sufficient for a readiness audit, but it is not a substitute for a proposal
engineering document.

The matrix also depends on current public status information. Unicode roadmaps,
pipeline pages, and proposal records can change over time
\cite{unicodeRoadmaps,unicodePipeline2026}. Any future use of the matrix for
standards-facing work should re-check public script-status sources. Similarly,
rights review remains necessary before any source image, glyph crop, or
manuscript example is distributed or reproduced.

Another limitation is that the present study cannot adjudicate the full history
of Nwagu Aneke use. It can cite the source anchor and public metadata, but it
does not claim a complete social history, complete manuscript census, or
complete account of contemporary community interest. Those questions are
important, but they require methods beyond a text-only readiness matrix:
archival work, oral history, authority review, and community-led decisions about
public representation.

Finally, the matrix does not replace community review. Community authority
cannot be inferred from the existence of a table, a citation, or a digital
model. It requires consultation, accountable representation, and decisions about
what should be made public. The matrix identifies that requirement; it does not
complete it.

These limitations constrain the strongest responsible claim. The paper can
state that Nwagu Aneke has enough evidence for a readiness audit and that the
current distribution of requirements blocks proposal framing. It cannot state
that the script is ready for encoding, that the character repertoire is
complete, that representative glyphs are settled, or that public release of
examples is authorized. Keeping those claims out of scope is not a rhetorical
weakness. It is the condition that makes the reported result reliable.

\section{Reproducibility}
The reported matrix can be reproduced by checking the twelve rows in
Table~\ref{tab:readiness} against the cited source and standards materials. The
minimum checks are: confirm that Nwagu Aneke is identifiable in Azuonye and
public script metadata; confirm that the working source layer is treated as 26
by 8 rather than 27/216; confirm that the f/v expansion is described as derived;
confirm that representative glyphs, usage examples, and community authority are
marked as blocked; and confirm that encoding model, punctuation and numbers,
implementation evidence, and proposal dossier are marked as missing.

Reproduction also requires checking the limitations. A reader should be able to
verify that source transcription review and rights review remain necessary
before stronger public claims are made. If new source records, rights decisions,
or community review outcomes become available, the relevant rows should be
updated rather than silently folded into the existing result.

For a second reviewer, the most important reproducibility check is the
source/derived distinction. The 26 by 8 table should reproduce as 208
source-observed records. The 27/216 figure should appear only as a derived f/v
expansion. If a future dataset, figure, or prose revision treats 27/216 as
source-observed, it has changed the evidentiary layer and should be rejected or
reclassified. This check is mechanical, but it protects the central claim of
the paper.

A third reviewer can reproduce the standards-facing conclusion by reading the
missing and blocked rows against Unicode proposal expectations. Usage,
stability, interchange need, character-glyph distinction, proposal documents,
character properties, examples, fonts, and licensing all appear in the relevant
standards guidance \cite{unicodeProposals,unicode17CoreSpec}. Under the present
matrix, several of those items are blocked or missing. The conclusion therefore
does not depend on a subjective judgment that the script is underdocumented; it
depends on row-level evidence that named readiness conditions have not yet been
satisfied.

\section{Conclusion}
Nwagu Aneke is ready for a Unicode-readiness audit, but the present evidence
does not support proposal framing. The twelve-item matrix finds one present
requirement, four partial requirements, three blocked requirements, and four
missing requirements. This result gives digital-humanities and standards
researchers a disciplined way to discuss an underdocumented script without
turning source promise into character-encoding readiness. The next scholarly
work is not a code chart; it is source transcription review, rights review,
community authority review, and character-evidence collection.

\bibliographystyle{ACM-Reference-Format}
\bibliography{../references}
\end{document}
